\documentclass[11pt,a4paper]{article}
\usepackage[T1]{fontenc}
\usepackage[utf8]{inputenc}
\usepackage{lmodern}
\usepackage{amsmath,amssymb,amsthm}
\usepackage{booktabs}
\usepackage{xcolor}
\usepackage[margin=27mm]{geometry}
\usepackage[colorlinks=true,urlcolor=blue,citecolor=blue,linkcolor=blue]{hyperref}
\newtheorem{theorem}{Theorem}
\newtheorem{proposition}[theorem]{Proposition}
\newtheorem{lemma}[theorem]{Lemma}
\newtheorem{corollary}[theorem]{Corollary}
\theoremstyle{definition}
\newtheorem{definition}[theorem]{Definition}
\newtheorem{hypothesis}[theorem]{Hypothesis}
\theoremstyle{plain}
\newcommand{\R}{\mathbb{R}}
\newcommand{\C}{\mathbb{C}}
\newcommand{\aTH}{\alpha_{\mathrm{TH}}}
\setlength{\emergencystretch}{2em}
\setlength{\parindent}{0pt}
\setlength{\parskip}{0.3em plus 0.1em minus 0.1em}
\title{Momentum R\'enyi and Tsallis entropies behind hard walls:\\
the existence threshold $1/(2(m+1))$}
\author{Fable \& Astra\\[0.35em]
\normalsize Research led by Jean-Philippe Bruneton\\[-0.1em]
\small Universit\'e Paris Cit\'e, France}
\date{19 November 2026}

\begin{document}
\maketitle
\begin{abstract}
A particle confined by impenetrable walls has a wave function that vanishes
outside its well. Cutting the wave function off at a wall produces an
algebraic tail in the momentum distribution, and R\'enyi or Tsallis entropies
of order $\alpha<1$, which weigh tails heavily, exist only when $\alpha$
exceeds a threshold $\aTH$. Olendski found $\aTH=1/4$ for Dirichlet walls and
$\aTH=1/2$ for Neumann and Robin walls in the one-dimensional well, and
conjectured that $1/2$ is the threshold of every one-dimensional
non-Dirichlet system. We prove a general criterion: if the first derivative
of the zero-extended wave function that jumps has order $m$, the momentum
entropies exist exactly for $\alpha>\aTH=1/(2(m+1))$, and the R\'enyi entropy
diverges logarithmically at the threshold with an explicit rate. For
eigenfunctions of a Schr\"odinger operator with integrable potential between
hard walls only $m=0$ and $m=1$ occur, so $\aTH=1/4$ when both walls are
Dirichlet and $\aTH=1/2$ otherwise. This proves Olendski's conjecture for
every level, for real or complex Robin lengths, and for any integrable
potential inside the well; the same holds on a half-line.
\end{abstract}

\section{Introduction}
A quantum particle in a well with impenetrable walls has a wave function that
is zero beyond the walls. Its momentum distribution is the squared modulus of
the Fourier transform of this truncated function, and the truncation costs a
slowly decaying tail at large momentum. Information-theoretic measures of the
momentum distribution feel this tail. For a normalised $\psi\in L^2(\R)$ let
\[
\phi(p)=\frac{1}{\sqrt{2\pi}}\int_\R\psi(x)e^{-ipx}\,dx,\qquad
\gamma(p)=|\phi(p)|^2,\qquad I_\alpha=\int_\R\gamma(p)^\alpha\,dp .
\]
The momentum R\'enyi and Tsallis \cite{Tsallis1988} entropies are
\[
R_\gamma(\alpha)=\frac{1}{1-\alpha}\ln I_\alpha,\qquad
T_\gamma(\alpha)=\frac{1}{\alpha-1}\bigl(1-I_\alpha\bigr),
\]
and for $\alpha\neq1$ they exist exactly when $I_\alpha<\infty$. When
$\alpha<1$ the power $\gamma^\alpha$ enhances the tail, so the entropies exist
only above a threshold $\aTH$ fixed by its decay rate.

Olendski computed these entropies for the one-dimensional well of width $a$
\cite{Olendski2020}. From the explicit momentum wave functions he found
$\aTH=1/4$ for the Dirichlet condition $\psi=0$ at the walls and $\aTH=1/2$
for the Neumann condition $\psi'=0$. Noting that the same value $1/2$ appears
for the attractive Robin wall and for mixed Dirichlet--Neumann walls, he
conjectured that ``the threshold $\aTH$ of one half is characteristic of any
1D non-Dirichlet system'' \cite{Olendski2020}. The class of systems is not
specified and no proof is given.

The simplest example shows the mechanism. The Neumann ground state
$\psi=a^{-1/2}$ on $[0,a]$ jumps from $0$ to $a^{-1/2}$ at each wall, so
$|\phi(p)|$ decays like $1/|p|$ and $\gamma^\alpha\sim|p|^{-2\alpha}$ is
integrable exactly for $\alpha>1/2$. The Dirichlet ground state is continuous,
its derivative jumps, $|\phi|$ decays like $|p|^{-2}$, and the threshold is
$1/4$. The general statement measures the order $m$ of the first derivative
that jumps.

Theorem~\ref{thm:threshold} shows that, if the zero-extended wave function has
its first jump in the derivative of order $m$, then $I_\alpha<\infty$ if and
only if $\alpha>1/(2(m+1))$. The delicate direction is divergence at and below
the threshold: with two walls, the leading coefficient of the momentum tail
is a trigonometric sum that may vanish periodically, and Lemma~\ref{lem:B}
shows that it stays large on a set of momenta of positive density.
Proposition~\ref{prop:rate} gives the exact divergence rate of $I_\alpha$ as
$\alpha\downarrow\aTH$. Corollary~\ref{cor:well} applies the theorem to
eigenfunctions of Schr\"odinger operators with integrable potential and
proves Olendski's conjecture in that setting.

\section{Related work}\label{sec:related}
Olendski's paper \cite{Olendski2020} determines the thresholds $1/4$ and $1/2$
for the Dirichlet and Neumann wells by a comparison test on explicit formulas,
and states the conjecture. Its introduction writes that the momentum entropies
exist ``at $\alpha\geq1/2$ only'', while its section on the Neumann well uses a
strict inequality; Theorem~\ref{thm:threshold} shows that the strict
inequality is correct. The attractive Robin wall had been treated earlier in
\cite{Olendski2019}. The author later extended the observation case by case to
other geometries: for two-dimensional circular dots he obtained $2/5$
(Dirichlet) and $2/3$ (Neumann) \cite{Olendski2021a}; for $d$-dimensional
balls, $d/(d+3)$ and $d/(d+1)$ \cite{Olendski2021b}, which reduce to $1/4$ and
$1/2$ for $d=1$; and for the pseudo-harmonic oscillator on the half-line, with
a potential containing a term $\mathfrak a/x^2$, a level-independent threshold
$1/(3+\sqrt{1+4\mathfrak a})$ \cite{Olendski2023}. All these thresholds are
obtained from explicit wave functions. Section~\ref{sec:scope} explains how the
ball and pseudo-harmonic values fit the formula $1/(2(m+1))$ heuristically.

The mechanism behind Lemma~\ref{lem:A} is classical: the large-frequency
asymptotics of a Fourier integral over an interval is governed by the jumps of
derivatives at the endpoints, as follows from repeated integration by parts;
see Erd\'elyi \cite{Erdelyi1955} for the general theory of such expansions.
What the present note adds is the full translation into an existence
criterion for the momentum entropies, including the divergence direction,
which requires controlling the interference between the contributions of two
walls, and the divergence rate at the threshold.

\section{Setting and statement}
\begin{hypothesis}[$H_m$]
Let $m\ge0$ be an integer. The function $\psi$ belongs to
$L^1(\R)\cap L^2(\R)$. If $m\ge1$, $\psi$ is of class $C^{m-1}$ on $\R$,
$\psi^{(m-1)}$ is absolutely continuous, and $\psi',\dots,\psi^{(m)}$ are
integrable on $\R$. There are points $x_1<\dots<x_N$, $N\ge1$, such that on
each open interval of $\R\setminus\{x_1,\dots,x_N\}$ the function
$\psi^{(m)}$ is locally absolutely continuous, and its derivative
$h=\psi^{(m+1)}$, defined almost everywhere, is integrable on $\R$. The jumps
$J_k=\psi^{(m)}(x_k^+)-\psi^{(m)}(x_k^-)$ are all nonzero.
\end{hypothesis}

The one-sided limits exist because $h$ is integrable, and for the same reason,
together with the integrability of $\psi^{(m)}$, $\psi^{(m)}$ tends to zero at
$\pm\infty$. The integer $m$ is the order of the first derivative of $\psi$
that jumps. For a smooth function on $[a,b]$ extended by zero, $N\le2$, the
points are the walls, $m$ is the smallest order such that $\psi^{(m)}(a)$ or
$\psi^{(m)}(b)$ is nonzero, and the jumps are $\psi^{(m)}(a)$ at $a$ and
$-\psi^{(m)}(b)$ at $b$ (only the nonzero ones are kept). Set
\[
G(p)=\sum_{k=1}^NJ_ke^{-ipx_k},\qquad \sigma^2=\sum_k|J_k|^2,\qquad
S=\sum_k|J_k| .
\]

\begin{lemma}[leading term]\label{lem:A}
Under $(H_m)$, for every $p\neq0$,
\[
\sqrt{2\pi}\,(ip)^{m+1}\phi(p)=G(p)+\hat h(p),\qquad
\hat h(p)=\int_\R h(x)e^{-ipx}\,dx\xrightarrow[|p|\to\infty]{}0 .
\]
\end{lemma}

\begin{proof}
For $j\le m$, $j$ integrations by parts on $\R$ give
$\int_\R\psi^{(j)}e^{-ipx}dx=(ip)^j\sqrt{2\pi}\,\phi(p)$; there are no
boundary terms because $\psi,\dots,\psi^{(m-1)}$ are absolutely continuous,
integrable with integrable derivatives, and hence tend to zero at infinity. On
each interval of $\R\setminus\{x_k\}$ one more integration by parts gives
\[
\int\psi^{(m)}e^{-ipx}dx=\Bigl[\frac{\psi^{(m)}e^{-ipx}}{-ip}\Bigr]
+\frac{1}{ip}\int he^{-ipx}dx .
\]
Summing over the intervals, the boundary terms at $x_k$ combine into
$J_ke^{-ipx_k}/(ip)$ and those at infinity vanish, so
$(ip)^m\sqrt{2\pi}\,\phi(p)=\bigl(G(p)+\hat h(p)\bigr)/(ip)$. The
Riemann--Lebesgue lemma gives $\hat h(p)\to0$.
\end{proof}

\begin{lemma}[the leading coefficient is large on a set of positive density]\label{lem:B}
If $N\ge2$, let $\delta=\min_{k\ne l}|x_k-x_l|$, $T_0=4S^2/(\sigma^2\delta)$
and $\mu_0=1/\delta$. For every real $P$, the set
$\Omega_P=\{p\in[P,P+T_0]:|G(p)|\ge\sigma/2\}$ has Lebesgue measure at least
$\mu_0$. If $N=1$, $|G|$ is constant and equal to $\sigma$, and the statement
holds with $T_0=\mu_0=1$.
\end{lemma}

\begin{proof}
Expanding the square,
\[
\frac1{T_0}\int_P^{P+T_0}|G|^2dp=\sigma^2+\sum_{k\ne l}J_k\bar J_l\,
\frac1{T_0}\int_P^{P+T_0}e^{-ip(x_k-x_l)}dp .
\]
Each oscillatory integral is bounded in modulus by $2/(T_0\delta)$, so the
cross terms are bounded by $2S^2/(T_0\delta)=\sigma^2/2$, and
$\int_P^{P+T_0}|G|^2dp\ge T_0\sigma^2/2$. Since $|G|\le S$ everywhere and
$|G|<\sigma/2$ outside $\Omega_P$, also
$\int_P^{P+T_0}|G|^2dp\le S^2|\Omega_P|+T_0\sigma^2/4$. Hence
$|\Omega_P|\ge T_0\sigma^2/(4S^2)=1/\delta$.
\end{proof}

Lemma~\ref{lem:B} is where the two walls interfere. The coefficient $G$ can
vanish: $G(p)=\psi(a)e^{-ipa}-\psi(b)e^{-ipb}$ has periodic zeros when
$|\psi(a)|=|\psi(b)|$. It cannot, however, be small on a proportion of
momenta large enough to restore convergence.

\begin{theorem}[threshold]\label{thm:threshold}
Under $(H_m)$ and for $\alpha>0$,
\[
I_\alpha=\int_\R\gamma(p)^\alpha\,dp<\infty
\quad\Longleftrightarrow\quad
\alpha>\aTH=\frac{1}{2(m+1)} .
\]
The momentum entropies $R_\gamma(\alpha)$ and $T_\gamma(\alpha)$ therefore
exist for $\alpha>\aTH$, $\alpha\neq1$, and do not exist for
$0<\alpha\le\aTH$.
\end{theorem}

\begin{proof}
The function $\phi$ is continuous and bounded by $\|\psi\|_1/\sqrt{2\pi}$, so
$\gamma^\alpha$ is locally integrable. By Lemma~\ref{lem:A} there is
$P_1\ge1$ such that $|\hat h(p)|\le\sigma/4$ for $|p|\ge P_1$. For such $p$,
$\gamma(p)\le(S+\sigma/4)^2/(2\pi|p|^{2m+2})$, so $I_\alpha<\infty$ as soon
as $2\alpha(m+1)>1$. Conversely, for $P\ge P_1$ and $p\in\Omega_P$ one has
$|G(p)+\hat h(p)|\ge\sigma/4$, hence $\gamma(p)\ge\sigma^2/(32\pi p^{2m+2})$
and
\[
\int_P^{P+T_0}\gamma^\alpha\,dp\ge\mu_0\Bigl(\frac{\sigma^2}{32\pi}\Bigr)^\alpha(P+T_0)^{-2\alpha(m+1)} .
\]
Summing over $P=P_1+jT_0$, $j\ge0$, gives a series whose general term is
equivalent to a constant times $j^{-2\alpha(m+1)}$, which diverges when
$2\alpha(m+1)\le1$.
\end{proof}

\begin{proposition}[divergence rate]\label{prop:rate}
Assume $(H_m)$ with $N\le2$, and let $M$ be the mean value of the periodic
function $|G|^{1/(m+1)}$: $M=|J_1|^{1/(m+1)}$ if $N=1$, and
\[
M=\frac1{2\pi}\int_0^{2\pi}\bigl(|J_1|^2+|J_2|^2+2|J_1||J_2|\cos\theta\bigr)^{\frac{1}{2(m+1)}}d\theta
\]
if $N=2$. Then
\[
\lim_{\alpha\downarrow\aTH}(\alpha-\aTH)\,I_\alpha=\frac{(2\pi)^{-\aTH}M}{m+1}>0,
\qquad
R_\gamma(\alpha)=\frac{1}{1-\aTH}\ln\frac{1}{\alpha-\aTH}+O(1) .
\]
\end{proposition}

\begin{proof}
We use an Abelian fact: if $f:[P,\infty)\to[0,U]$ is measurable, $P\ge1$, and
$\frac1T\int_P^{P+T}f\to M$ as $T\to\infty$, then
$(s-1)\int_P^\infty f(p)p^{-s}dp\to M$ as $s\downarrow1$. Indeed, with
$F(p)=\int_P^pf\le Up$, an integration by parts gives
$\int_P^\infty fp^{-s}dp=s\int_P^\infty F(p)p^{-s-1}dp$. For $\varepsilon>0$
there is $P_\varepsilon$ with $|F(p)-Mp|\le\varepsilon p$ for $p\ge P_\varepsilon$,
whence
\[
\Bigl|s\int_P^\infty Fp^{-s-1}dp-sM\int_P^\infty p^{-s}dp\Bigr|
\le s(U+M)P_\varepsilon+\frac{s\varepsilon}{s-1};
\]
multiplying by $s-1$ and letting $s\to1$, then $\varepsilon\to0$, proves the
claim.

Set $s=2\alpha(m+1)$, so $s-1=2(m+1)(\alpha-\aTH)$. For $p\ge P_1$,
Lemma~\ref{lem:A} reads $\gamma(p)^\alpha=(2\pi)^{-\alpha}p^{-s}f_\alpha(p)$
with $f_\alpha=|G+\hat h|^{2\alpha}$. Since
$\bigl||G+\hat h|-|G|\bigr|\le|\hat h|\to0$ and $(u,\alpha)\mapsto u^{2\alpha}$
is uniformly continuous on $[0,2S]\times[\aTH,1]$, for every $\eta>0$ there
are $P_\eta\ge P_1$ and $\alpha_\eta>\aTH$ such that
$|f_\alpha(p)-|G(p)|^{2\aTH}|\le\eta$ for $p\ge P_\eta$ and
$\aTH<\alpha\le\alpha_\eta$; the corresponding error on
$(s-1)\int_{P_\eta}^\infty f_\alpha p^{-s}dp$ is at most
$\eta P_\eta^{1-s}\le\eta$. For $N\le2$, $|G|^{2\aTH}=|G|^{1/(m+1)}$ is
periodic (constant if $N=1$) with mean $M$, so the Abelian fact gives the
limit $M$ on the positive half-line. The same holds on the negative half-line,
since $|G(-p)|^2=|J_1|^2+|J_2|^2+2|J_1||J_2|\cos(\cdot)$ has the same mean. The
contribution of $|p|\le P_\eta$ is bounded uniformly in $\alpha$ and
disappears after multiplication by $s-1$. Altogether
$(s-1)I_\alpha\to2(2\pi)^{-\aTH}M$, which is the stated limit; $M>0$ by
Lemma~\ref{lem:B}. The expansion of $R_\gamma$ follows by taking logarithms.
\end{proof}

The divergence of the R\'enyi entropy at the threshold is therefore
logarithmic, as Olendski's figures show \cite{Olendski2020}, with the
universal coefficient $1/(1-\aTH)$: $2$ for a non-Dirichlet wall and $4/3$ for
the Dirichlet well. For the Neumann ground state of the well of width $a$,
$\psi=a^{-1/2}$, $J_1=-J_2=a^{-1/2}$, $M=4/(\pi\sqrt a)$, and
\[
\lim_{\alpha\downarrow1/2}\Bigl(\alpha-\frac12\Bigr)I_\alpha
=\frac{4}{\pi\sqrt{2\pi a}}\approx\frac{0.5079}{\sqrt a}.
\]
For $N\ge3$ the proposition should remain true with $M$ the Bohr mean value of
the almost periodic function $|G|^{1/(m+1)}$; this extension is not written
out here.

\section{Hard-wall eigenfunctions: the conjecture}
\begin{corollary}[well with hard walls]\label{cor:well}
Let $a<b$, let $V\in L^1(a,b)$ be a complex-valued potential and $E\in\C$, and
let $\psi$ be a solution of $-\psi''+V\psi=E\psi$ on $[a,b]$, not identically
zero, with $\psi$ and $\psi'$ absolutely continuous, extended by zero outside
$[a,b]$. If $\psi(a)=\psi(b)=0$, then $\aTH=1/4$; in every other case
$\aTH=1/2$. In particular, if one of the walls carries a Robin condition
$\psi'(a)=\psi(a)/\Lambda_a$ or $-\psi'(b)=\psi(b)/\Lambda_b$ with
$\Lambda\in\C\setminus\{0\}$, or a Neumann condition ($\Lambda=\infty$), then
every eigenfunction, at every level, has $\aTH=1/2$.
\end{corollary}

\begin{proof}
By the uniqueness theorem for linear differential equations with integrable
coefficients, a solution with $\psi(c)=\psi'(c)=0$ at some $c\in[a,b]$
vanishes identically. If $\psi(a)\neq0$ or $\psi(b)\neq0$, the extended
function satisfies $(H_0)$: it is bounded with compact support, its derivative
$h=\psi'$ on $(a,b)$ is integrable, and its jumps are $\psi(a)$ at $a$ and
$-\psi(b)$ at $b$, of which only the nonzero ones are kept; where $\psi$
vanishes at a wall the extension is continuous there. If $\psi(a)=\psi(b)=0$,
the extended function is absolutely continuous on $\R$, its derivative jumps
by $\psi'(a)\ne0$ at $a$ and by $-\psi'(b)\neq0$ at $b$ by uniqueness, and
$h=\psi''=(V-E)\psi$ is integrable on $(a,b)$ because $\psi$ is bounded: this
is $(H_1)$. Theorem~\ref{thm:threshold} concludes. Finally, at a Robin or
Neumann wall, $\psi(a)=0$ would force $\psi'(a)=0$ and hence $\psi\equiv0$, so
$\psi(a)\neq0$.
\end{proof}

\begin{corollary}[half-line]\label{cor:half}
Let $\psi$ be a nonzero solution of the same equation on $[0,\infty)$, with
$V$ locally integrable on $[0,\infty)$, extended by zero for $x<0$, such that
$\psi$, $\psi'$ and $\psi''$ are integrable on $(0,\infty)$. If
$\psi(0)\neq0$, then $\aTH=1/2$; if $\psi(0)=0$ and $\psi'(0)\neq0$, then
$\aTH=1/4$.
\end{corollary}

\begin{proof}
Hypotheses $(H_0)$ and $(H_1)$ are verified as in Corollary~\ref{cor:well},
with $N=1$; $\psi\in L^2$ because $\psi$ is integrable and bounded.
\end{proof}

The attractive Robin wall, with bound state $\psi(x)=\sqrt{2\kappa}\,e^{-\kappa x}$
on the half-line \cite{Olendski2019}, falls under the first case. Here
$\gamma(p)=\kappa/(\pi(\kappa^2+p^2))$ and
\[
I_\alpha=\Bigl(\frac\kappa\pi\Bigr)^\alpha\kappa^{1-2\alpha}\sqrt\pi\,
\frac{\Gamma(\alpha-\frac12)}{\Gamma(\alpha)},
\]
finite exactly for $\alpha>1/2$, with a simple pole of residue
$\sqrt{\kappa/\pi}=(2\pi)^{-1/2}|J_1|$, as Proposition~\ref{prop:rate}
predicts. The quasi-one-dimensional hydrogen ground state, $\psi(x)\propto
xe^{-x}$, satisfies $(H_1)$ directly with a single jump of $\psi'$ at the
origin and recovers the threshold $1/4$ reported by Olendski.

The theorem also explains why only two values occur in these models: an
eigenfunction of a Schr\"odinger operator with integrable potential cannot
have both $\psi$ and $\psi'$ zero at a wall, so $m\in\{0,1\}$. Values $m\ge2$
require functions that are not such eigenfunctions, for example
$\psi\propto x^2(1-x)^2$ on $[0,1]$, whose second derivative jumps by $2$ at
$0$ and by $-2$ at $1$, so that $\aTH=1/6$.

\section{Verification}\label{sec:verification}
The statements above are proved analytically. The public
\href{https://github.com/jpbruneton/open_quantum_problems/blob/main/research/entries/momentum-renyi-threshold/checks/check_momentum_threshold.py}{Python checker}
illustrates them numerically and writes its
\href{https://github.com/jpbruneton/open_quantum_problems/blob/main/research/entries/momentum-renyi-threshold/data/verification.json}{verification record}
to the same entry. It builds nine eigenfunctions of the well $[0,1]$ with
Neumann, Dirichlet, mixed, real Robin ($\Lambda=0.5$, $-0.3$, $2$, $0.7$) and
complex Robin ($\Lambda=0.4+0.3i$) walls, whose Fourier transforms are known
in closed form and are checked against quadrature to about $10^{-16}$. It
verifies that $m=1$ exactly when both walls are Dirichlet; that the difference
between $|p|^{m+1}\sqrt{2\pi}\,|\phi(p)|$ and $|G(p)|$, multiplied by $P$,
stays bounded on windows at $\pm P$ for $P=10^2,10^3,10^4$
(Lemma~\ref{lem:A}); that the dyadic blocks $\int_P^{2P}\gamma^{\aTH}dp$
approach the nonzero constant $\ln2\,(2\pi)^{-\aTH}M$, within a relative
$2\times10^{-3}$ at $P=6400$, which is the divergence at the threshold; and
that the ratio of successive blocks equals $2^{1-2\alpha(m+1)}$ within
$5\times10^{-3}$ for $\alpha=\aTH\pm0.1$. It also treats the polynomial example
with $m=2$, the closed form for the attractive Robin wall, and the divergence
rate of the Neumann ground state through a Hurwitz zeta representation:
$(\alpha-\frac12)I_\alpha=0.5131$, $0.5085$ and $0.5080$ for
$\alpha-\frac12=10^{-2},10^{-3},10^{-4}$, against the limit $0.5079$. The
checker requires NumPy, SciPy and mpmath and runs in under a minute. These
computations illustrate the proof and do not replace it.

\section{Scope and limitations}\label{sec:scope}
\emph{Analytic proof.} Lemmas~\ref{lem:A} and~\ref{lem:B},
Theorem~\ref{thm:threshold}, Proposition~\ref{prop:rate} for $N\le2$ and
Corollaries~\ref{cor:well} and~\ref{cor:half}. The conjecture is thus
established for a particle on a bounded interval or a half-line, with a wave
function vanishing beyond hard walls and solving a Schr\"odinger equation with
integrable potential (locally integrable near the wall on the half-line), with
at least one wall where the condition is not Dirichlet. It holds for every
level, for real or complex Robin lengths, and also for superpositions and
non-stationary states whenever $\psi$ does not vanish at a wall and satisfies
$(H_0)$.

\emph{Not covered.} A one-dimensional system without hard walls has no
threshold of this type, and ``non-Dirichlet'' has no meaning there. Potentials
that are singular at the wall, such as $\mathfrak a/x^2$, produce
$\psi\sim x^\nu$ with non-integer $\nu$. Olendski's value
$1/(3+\sqrt{1+4\mathfrak a})$ for the pseudo-harmonic oscillator
\cite{Olendski2023} coincides with $1/(2(\nu+1))$ for
$\nu=\frac12(1+\sqrt{1+4\mathfrak a})$, the natural extension of
Theorem~\ref{thm:threshold}, but this fractional case is not proved here. In
dimension $d$, the thresholds $d/(d+1)$ and $d/(d+3)$ found for balls
\cite{Olendski2021b} are those given by the heuristic
$\gamma\sim|p|^{-(d+1+2m)}$ integrated against $|p|^{d-1}d|p|$, namely
$\aTH=d/(d+1+2m)$; no general proof is attempted for arbitrary domains.

\emph{Bibliographic priority.} The source \cite{Olendski2020} was read in
full. The twelve works citing it listed by OpenAlex (2020 to 2026) were
examined by title only; the author's later preprints were examined through
their abstracts, and the relevant passage of \cite{Olendski2023} in its
source. Two web searches on the existence of momentum R\'enyi entropies, boundary
conditions and Fourier decay returned no general statement. The asymptotic
expansion of Lemma~\ref{lem:A} is classical, and Olendski already asserts the
threshold $1/2$ for the Robin well on the basis of a comparison test. The
general statement in terms of $m$, the proof of divergence with two walls and
the divergence rate were not found in this documented search, which does not
establish priority.

\clearpage
\begingroup\small
\begin{thebibliography}{99}
\setlength{\itemsep}{2pt}
\bibitem{Olendski2020}
O. Olendski,
\emph{R\'enyi and Tsallis entropies of the Dirichlet and Neumann
one-dimensional quantum wells},
Int. J. Quantum Chem. \textbf{120}, e26220 (2020);
\href{https://doi.org/10.1002/qua.26220}{doi:10.1002/qua.26220};
\href{https://arxiv.org/abs/2003.04351}{arXiv:2003.04351}.

\bibitem{Olendski2019}
O. Olendski,
\emph{R\'enyi and Tsallis entropies: three analytic examples},
Eur. J. Phys. \textbf{40}, 025402 (2019);
\href{https://doi.org/10.1088/1361-6404/aaf45c}{doi:10.1088/1361-6404/aaf45c};
\href{https://arxiv.org/abs/1811.09089}{arXiv:1811.09089}.

\bibitem{Olendski2021a}
O. Olendski,
\emph{Comparative analysis of information measures of the Dirichlet and
Neumann two-dimensional quantum dots},
Int. J. Quantum Chem. \textbf{121}, e26455 (2021);
\href{https://doi.org/10.1002/qua.26455}{doi:10.1002/qua.26455};
\href{https://arxiv.org/abs/1806.07509}{arXiv:1806.07509}.

\bibitem{Olendski2021b}
O. Olendski,
\emph{Quantum information measures of the Dirichlet and Neumann hyperspherical
dots},
Eur. Phys. J. Plus \textbf{136}, 390 (2021);
\href{https://doi.org/10.1140/epjp/s13360-021-01366-3}{doi:10.1140/epjp/s13360-021-01366-3};
\href{https://arxiv.org/abs/2103.13094}{arXiv:2103.13094}.

\bibitem{Olendski2023}
O. Olendski,
\emph{One-dimensional pseudoharmonic oscillator: classical remarks and
quantum-information theory}, preprint (2023);
\href{https://arxiv.org/abs/2304.06428}{arXiv:2304.06428}.

\bibitem{Erdelyi1955}
A. Erd\'elyi,
\emph{Asymptotic Representations of Fourier Integrals and the Method of
Stationary Phase},
J. Soc. Indust. Appl. Math. \textbf{3}, 17--27 (1955);
\href{https://doi.org/10.1137/0103002}{doi:10.1137/0103002}.

\bibitem{Tsallis1988}
C. Tsallis,
\emph{Possible generalization of Boltzmann-Gibbs statistics},
J. Stat. Phys. \textbf{52}, 479--487 (1988);
\href{https://doi.org/10.1007/BF01016429}{doi:10.1007/BF01016429}.
\end{thebibliography}
\endgroup
\end{document}
